% Pro M. Caelio --- headnote
% pro-caelio, 4 April 56 BC, Rome

\textit{Cicero for M. Caelius Rufus, delivered at Rome on 4
April 56 BC --- the second day of the Megalensian games, on
which the courts were ordinarily not in session. The
prosecution had nevertheless got the trial set for the
festival, on the law on \textit{vis} (the Plautian law on
public violence), so that Caelius's youthful brilliance, his
recent prosecution of L. Bestia (twice, the second time after
acquittal), and the inflated catalogue of charges around the
Egyptian commission of Ptolemy Auletes might be heard before a
half-empty court. The defence team was Caelius himself
(speaking for himself, on the second day), then M. Crassus, then
Cicero closing. The prosecutors were L. Sempronius Atratinus
(son of L. Calpurnius Bestia, whom Caelius was prosecuting), L.
Herennius Balbus, and P. Clodius (the son of Appius Claudius
Pulcher, not the famous tribune --- a young man).}

\textit{The actual charges, as Cicero shapes them, are five
counts of \textit{vis}: \textsection 23 the Neapolitan
seditions, the assault on the Alexandrian envoys at Puteoli, the
business of Palla's property, the death of Dio of Alexandria;
\textsection 30 the gold from Clodia, and the poison alleged to
have been prepared for Clodia herself. The first three Cicero
turns over to Crassus's portion of the speech and refers
to as already disposed of. The Dio charge collapses on the
acquittal of P. Asicius (whom Cicero himself defended), the
confession of Ptolemy, and the testimony Cicero summons from
the Coponii brothers, Dio's hosts. The case as Cicero handles
it reduces to two charges --- gold and poison --- both of which
``one and the same person is involved'' (\textsection 30): Clodia.}

\textit{The speech is the most virtuosic forensic rhetoric
Cicero ever wrote. The ground is shifted from criminal trial
to social comedy: Clodia, the elder sister of P. Clodius
Pulcher and widow of Q. Metellus Celer (consul of 60 BC, dead
in late 59 BC), is the prosecution's witness and source of the
charges, and Cicero turns her into the case. The interlocking
prosopopoeiae of \textsection\textsection 33--38 are the
centerpiece: Appius Claudius Caecus the censor, Clodia's blind
ancestor, summoned from the underworld in Republican gravitas
to scold her (\textsection\textsection 33--34), with the closing
catalogue of his works (the broken peace with Pyrrhus, the Aqua
Appia, the Via Appia) turned bitter --- ``did I lay down a road
that you should throng it with strangers' husbands?''; her
``littlest brother'' (Clodius) in urbane mode advising her on
courtesan economy (\textsection 36); the Caecilian comic
``hard father'' and the merciful father, in succession
(\textsection\textsection 37--38). Around them are the great set
pieces: the Luperci joke (\textsection 26 ``a wild brotherhood
of the German Luperci, quite of the shepherds and the country,
their woodland coming-together instituted before humanity and
the laws''); the Crassus-Medea allusion in the speech on the
Palatine apartment (\textsection 18 ``\textit{utinam ne in nemore
Pelio}\dots Medea sick at heart, by savage love wounded'');
Baiae crying out at \textsection 47; the long conditional on
the courtesan-and-lover at \textsection\textsection 48--49; and
the Senian-baths poison-trap unwound as a mime in
\textsection\textsection 61--65 --- ``the end then is of a mime,
not a play, in which when no \textit{clausula} is found
somebody slips out of the hands, then the clappers crash, the
curtain rises.''}

\textit{The Metellus death-bed digression at
\textsection\textsection 59--60 (Cicero's voice ``weakened with
weeping and the mind hindered with grief'' at remembering it)
plants the unspoken charge: that Clodia poisoned her own
husband. The structure of innuendo is exact: from this house
the speaker about poison's swiftness has come forth, will she
not fear the conscious walls and the funereal night? The
\textsection\textsection 75 admission of Caelius's ``turning''
of his age, paused at the goal-posts in the new-acquaintance
with Clodia and the unhappy neighbourhood, is the speech's
single concession --- and immediately framed as the broken-off
moment from which Caelius emerged with the prosecution of
Bestia.}

\textit{The peroration (\textsection\textsection 77--80) is the
classic two-prosecutions hostage formula --- a man who has
prosecuted a consular for violation of the commonwealth cannot
himself be a turbulent citizen, a man who twice prosecutes for
canvassing cannot himself be a briber --- and the close on
Caelius's old father, leaning on this only son: the
\textit{senex} of Roman comedy turned to peroration, with the
sharp closing aside on the recent acquittal of Sex. Clodius
(``a man without property, without faith, without hope, without
seat, without fortunes, polluted in face, tongue, hand, and his
whole life'') as the warning that, in the same court, with the
same woman, the foulest brigand should not be acquitted by a
woman's favour while the most honourable young man is delivered
up to a woman's lust.}

\textit{Caelius was acquitted. The speech secured for him the
opening of a brilliant if turbulent career: tribune of the plebs
in 52, praetor in 48 (a Caesarian who broke with Caesar over
debt-cancellation, killed in southern Italy in early 48 BC).
Clodia is not certainly heard from again in the historical
record. The speech survives as the great example of Roman
forensic rhetoric in the urbane register --- the trial-defence
turned into social comedy by sustained \textit{prosopopoeia} ---
and as one of the chief documents for the social world of the
late Republic at its dazzling and corroded surface. The
identification of Clodia with the ``Lesbia'' of Catullus's
poems is widely held but not unanimous (the metric pseudonym
suggests it; Apuleius asserts it; the chronology --- Catullus's
death conventionally placed ca. 54 BC --- is consistent).}
